\chapter{Informations, classement et sélection}
\section{Lire les informations}
Le panneau gauche décrit la photo ouverte : nom, date disponible, stockage,
dimensions, appareil et autres métadonnées connues. Une information absente
ne doit pas être interprétée comme une valeur mesurée égale à zéro.
Sous l’histogramme, les informations de prise de vue disponibles facilitent
la lecture des ISO, de la focale et du temps d’exposition.

\section{Noter et marquer une photo}
Cliquer sur une des cinq étoiles pour donner une note de 1 à 5.
Cliquer sur la note déjà sélectionnée la remet à zéro. Le drapeau se bascule
par un clic. Choisir un libellé rouge, jaune, vert, bleu ou violet ; cliquer
sur la même couleur retire le libellé. Ces choix de classement ne recolorent
pas les pixels et ne correspondent pas aux bandes du mélangeur.

Aucun raccourci numérique de notation ou de drapeau n’est déclaré pour
l’instant. Les commandes se trouvent dans le panneau Informations.

\section{Copier les retouches à plusieurs photos}
\begin{enumerate}
\item Ouvrir la photo portant les réglages à copier.
\item Cocher les cases des vignettes destinataires, sans confondre case cochée
et image ouverte. Le compteur indique le nombre de fichiers sélectionnés.
\item Cliquer sur \emph{Appliquer à N}. Vérifier le message de résultat.
\item Décocher les fichiers, ou utiliser le bouton de désélection générale.
\end{enumerate}
Cette action écrit les retouches dans les compagnons des destinataires.
Elle n’est pas limitée à la seule famille actuellement dépliée. Les profils
RAW et la balance automatique propres aux destinataires ne sont pas copiés
aveuglément entre boîtiers ; leurs métadonnées de classement sont conservées.

\attention{Il s’agit d’une application groupée des retouches, pas d’une commande
d’export multiple. Il n’existe pas de recette utilisateur de sélection par
Ctrl-clic ou Maj-clic dans le ruban actuel : utiliser les cases.}

\section{Les badges du ruban}
Le badge de la vignette ouverte peut rappeler sa note, son drapeau et son
libellé. L’absence de badge sur les autres vignettes ne signifie pas qu’elles
ne sont pas notées ou retouchées : ces données ne sont pas encore exposées
pour toutes les lignes du modèle.

\section{Clic droit sur une vignette}
Le menu vise la vignette sous la souris. Sur une photo différente de celle
ouverte, choisir d’abord \emph{Ouvrir cette image} : le clic droit seul ne
la charge pas et ne change pas la sélection multiple. Les actions de retouche
restent désactivées pour cette cible tant qu’elle n’est pas ouverte, afin
de ne pas réinitialiser ou exporter la photo précédente par erreur.

\emph{Dupliquer le fichier et ses réglages}, associé au raccourci réservé
\touche{Ctrl + D}, doit créer une copie physique et son compagnon cohérent,
pas un nouveau TIFF développé ni une copie virtuelle. Cette commande attend
encore l’interface moteur et reste désactivée dans la version décrite.
